\section{NormCode Canvas: The Deployed System}
\label{sec:canvas}

NormCode Canvas (v1.1.3) is the deployed system that makes both levels of CBR
accessible in one environment. The CBR operations are its primary interaction model.

\subsection{Level 1 Interface: Concrete Case Management}

\textbf{Graph View.} Before execution, Canvas renders the complete inference graph from
the compiled \texttt{concept\_repo.json} and \texttt{inference\_repo.json}. Every node
is a potential Level~1 case (a future checkpoint), identified by its flow index. Semantic
nodes (purple hexagons) and syntactic nodes (gray rectangles) are visually distinct.

\textbf{Real-Time Execution --- Retain.} During a run, Canvas streams execution events
via WebSocket. Every transition to completed writes a new Level~1 case to SQLite:
(\texttt{run\_id}, \texttt{flow\_index}, Blackboard snapshot, Concept Repository
snapshot). Retain is fully automatic.

\textbf{Breakpoint + Tensor Inspector --- Retrieve.} Any node can be designated as a
breakpoint by flow index. Execution pauses before that node runs, giving the operator
the opportunity to inspect the case. The Tensor Inspector displays the full
$N$-dimensional Reference: axis names, shapes, and content values or perceptual signs.
Inspection is O(1) --- C1.

\textbf{Value Override --- Revise.} At a breakpoint, or from the Checkpoint Panel on
any completed case, the operator clicks Override to enter edit mode. On resume, the
orchestrator records the modified tensor, marks all downstream nodes as stale (computed
automatically from the dependency graph), and re-executes only stale nodes --- C3.

\textbf{Fork --- Reuse.} The Checkpoint Panel lists all Level~1 cases organized by run
ID and flow index. Fork instantiates a new execution from any checkpoint: new run $r'$,
load stored $(\mathcal{B}_f, \mathcal{C}_f)$, all upstream marked completed, all
downstream execute fresh. Fork supports failure recovery, A/B testing, cost efficiency,
and iterative refinement.

\subsection{Level 2 Interface: Abstract Case Management}

\textbf{The \texttt{.ncn} Review Panel.} Canvas exposes the \texttt{.ncn} in a review
panel available before the first run. A domain expert reads the plain-prose narrative,
verifies reasoning structure matches intent, and either approves or rejects for revision.
This is C2 --- zero-cost pre-execution review.

\textbf{Compilation Workflow Visibility.} When the NC~Compilations plan runs the
compiler pipeline, Canvas shows it as a Level~1 execution: each compilation phase is a
node producing checkpoints. The compilation pipeline is thus debuggable via Level~1 CBR.

\textbf{Canvas Assistant.} Canvas ships a built-in NormCode plan --- the Canvas
Assistant --- that runs inside the Canvas App itself. Users interact through a chat panel;
the plan classifies each natural-language message, dispatches it to the appropriate
Canvas Integration facility (navigation via \texttt{me.hands}, execution control via
\texttt{me.mind}, or state query via \texttt{me.vision}), and responds through the same
chat. The plan is itself a live Level~1 CBR object: its per-iteration checkpoints are
inspectable and revisable in the Tensor Inspector. This is a concrete instance of a plan
using Canvas Integration facilities to control the very environment in which it runs.

\subsection{Architecture}

\textbf{Frontend:} React~18 with TypeScript; React Flow for graph visualization; Zustand
for state management; TailwindCSS.

\textbf{Backend:} FastAPI with Python~3.11; SQLite for case base persistence; native
Python WebSockets for real-time execution events.

\textbf{Orchestrator:} Dependency-driven scheduling via Waitlist (priority queue by flow
index) and Blackboard (status tracker). The scope rule is enforced by construction: the
orchestrator passes only the declared tensor references to each inference's input ---
never the full Concept Repository.

\textbf{Multi-agent support.} Multiple agents with different LLM models (e.g.,
\texttt{qwen-plus}, \texttt{gpt-4o}, \texttt{claude-3-opus}) can be registered.
Inferences are mapped to agents by pattern rules on flow indices.
